\documentclass[aspectratio=169,11pt]{beamer}

\usepackage[UTF8]{ctex}
\usepackage{tikz}
\usepackage{fontspec}
\usepackage{xcolor}
\usepackage{booktabs}
\usepackage{array}
\usepackage{graphicx}
\usepackage{hyperref}

\usetikzlibrary{arrows.meta,positioning,calc,shapes.geometric}

\setCJKmainfont{Microsoft YaHei}
\setsansfont{Arial}
\setmainfont{Arial}

\definecolor{Ink}{HTML}{17211D}
\definecolor{Moss}{HTML}{2D5A4B}
\definecolor{Sage}{HTML}{739A83}
\definecolor{Warm}{HTML}{FFFFFF}
\definecolor{Line}{HTML}{D7D0C4}
\definecolor{Clay}{HTML}{B16E45}
\definecolor{Mist}{HTML}{EBEFE8}
\definecolor{FootA}{HTML}{DCE5DC}
\definecolor{FootB}{HTML}{CAD9CC}
\definecolor{FootText}{HTML}{F4F1E8}

\setbeamertemplate{navigation symbols}{}
\setbeamertemplate{headline}{}
\setbeamertemplate{frametitle}{%
  \nointerlineskip
  \hspace*{-1.0cm}%
  
\begin{tikzpicture}
    \useasboundingbox (0,0) rectangle (\paperwidth,0.88cm);
    \fill[Moss] (0,0) rectangle (\paperwidth,0.88cm);
    \fill[Sage] (0,0) rectangle (0.26cm,0.88cm);
    \node[anchor=west,text=Warm,font=\bfseries\large] at (0.55cm,0.44cm) {\insertframetitle};
  \end{tikzpicture}
  \vspace{0.1cm}
}
\setbeamertemplate{footline}{%
  \nointerlineskip
  
\begin{tikzpicture}
    \useasboundingbox (0,0) rectangle (\paperwidth,0.46cm);
    \fill[FootA] (0,0) rectangle (\paperwidth,0.46cm);
    \fill[Moss] (0,0) rectangle (3.45cm,0.46cm);
    \fill[Moss] ($(0,0)+(\paperwidth-1.28cm,0)$) rectangle (\paperwidth,0.46cm);
    \node[anchor=west,text=FootText,font=\scriptsize] at (0.34cm,0.23cm) {MetaSim};
    \node[anchor=center,text=FootText,font=\scriptsize] at ($(0,0)+(\paperwidth-0.64cm,0.23cm)$) {\insertframenumber/\inserttotalframenumber};
  \end{tikzpicture}
}
\setbeamercolor{frametitlebar}{bg=Moss,fg=Warm}
\setbeamercolor{footupper}{bg=FootA}
\setbeamercolor{footlower}{bg=Moss}
\setbeamercolor{background canvas}{bg=Warm}
\setbeamercolor{normal text}{fg=Ink}
\setbeamercolor{frametitle}{fg=Ink}
\setbeamerfont{frametitle}{size=\Large,series=\bfseries}
\setbeamerfont{title}{size=\huge,series=\bfseries}
\setbeamerfont{subtitle}{size=\large}

\newcommand{\bigclaim}[1]{{\color{Ink}\fontsize{22}{27}\selectfont\bfseries #1}}
\newcommand{\keyclaim}[1]{{\color{Ink}\fontsize{26}{30}\selectfont\bfseries #1}}
\newcommand{\muted}[1]{{\color{Sage}#1}}
\newcommand{\pill}[1]{%
  \tikz[baseline=(x.base)] \node[inner xsep=7pt,inner ysep=3pt,rounded corners=2pt,fill=Mist,draw=Line,text=Ink,font=\scriptsize] (x) {#1};%
}
\newcommand{\metric}[2]{%
  \begin{minipage}{0.29\linewidth}
  {\color{Moss}\fontsize{22}{24}\selectfont\bfseries #1}\par
  \vspace{0.08cm}{\color{Ink}\small #2}
  \end{minipage}
}
\newcommand{\capblock}[2]{%
  \begin{minipage}{0.31\linewidth}
    {\color{Moss}\large\bfseries #1}\par
    \vspace{0.12cm}{\small #2}
  \end{minipage}
}

\title{MetaSim}
\subtitle{An Extensible LLM-based Generative Agent Simulation Platform and Applications in Crisis Management}
\date{}
\author{}
\institute{}

\begin{document}

{
\setbeamertemplate{footline}{}
\setbeamertemplate{headline}{}
\setbeamertemplate{frametitle}{}
\begin{frame}[plain]
  \begin{tikzpicture}[remember picture,overlay]
    \node[anchor=south west,inner sep=0pt] at (current page.south west) {
      
\includegraphics[width=\paperwidth,height=\paperheight]{slides/title.png}
    };
  \end{tikzpicture}
\end{frame}
}

\setbeamertemplate{headline}{}

\begin{frame}{定位}
    \keyclaim{Meta Simulator}

  \vspace{0.55cm}
  \begin{columns}[T,totalwidth=\linewidth]
    \begin{column}{0.48\linewidth}
      \small
      MetaSim 是一个基于 ECS（Entity-Component System）与数据驱动的离散事件仿真沙盒内核。提供世界建模、规则运行、状态结算、智能体接入和过程记录等基础能力。

      \vspace{0.35cm}
      开发者可以把不同主题的场景快速放入同一个运行框架中，例如野营、生存、狼人杀、竞技场、陪伴机器人或教学环境。
    \end{column}
    \begin{column}{0.46\linewidth}
      \vspace{-1.02cm}
      \centering
      \includegraphics[width=\linewidth,height=5.95cm,keepaspectratio,trim=0 0 0 1.15cm,clip]{slides/定位.png}
    \end{column}
  \end{columns}
\end{frame}

\begin{frame}{能力}
  \centering
  \vspace{-0.18cm}
  \includegraphics[width=0.96\paperwidth,height=6.05cm,keepaspectratio]{slides/能力.png}
\end{frame}

\begin{frame}{流程}
    \centering
    \vspace{-0.28cm}
    \includegraphics[width=\paperwidth,height=6.35cm,keepaspectratio]{slides/流程.png}
\end{frame}

\begin{frame}{边界}
    \centering
    \vspace{0.05cm}
    {\color{Ink}\small MetaSim 系统在大语言模型代理尝试不存在的动作时，防止动作幻觉，并引导其重新规划。}

    \vspace{0.08cm}
    \includegraphics[width=\paperwidth,height=4.95cm,keepaspectratio]{slides/边界.png}

    \vspace{0.08cm}
    {\color{Ink}\small 世界状态保持不变。只有经过验证的配方/效果才能修改世界。}
\end{frame}

\begin{frame}{迁移}
    \begin{columns}[T,totalwidth=0.92\paperwidth]
      \begin{column}{0.34\paperwidth}
        \small
        MetaSim 可以扩展到社交平台中的风险内容传播模拟。

        \vspace{0.22cm}
        独立 runtime 接入后，社交平台成为 agent 可进入的外部环境。

        \vspace{0.28cm}
        {\color{Sage}\footnotesize\bfseries 手机屏幕}\par
        {\footnotesize feed、post、操作反馈}

        \vspace{0.18cm}
        {\color{Sage}\footnotesize\bfseries 平台 runtime}\par
        {\footnotesize 账号、帖子、推荐、互动}

        \vspace{0.18cm}
        {\color{Sage}\footnotesize\bfseries MetaSim 记录}\par
        {\footnotesize 事件、记忆、指标}
      \end{column}
      \begin{column}{0.54\paperwidth}
        \centering
        \vspace{0.08cm}
        \includegraphics[width=\linewidth,height=5.02cm,keepaspectratio]{slides/迁移.png}
      \end{column}
    \end{columns}
\end{frame}

\begin{frame}{社交平台}
    \begin{columns}[T,totalwidth=0.92\paperwidth]
      \begin{column}{0.34\paperwidth}
        \small
        社交平台页展示 agent 与平台 runtime 的具体交互方式。

        \vspace{0.22cm}
        agent 只通过手机屏幕看到平台内容，并通过结构化动作影响平台状态。

        \vspace{0.28cm}
        {\color{Sage}\footnotesize\bfseries 可见内容}\par
        {\footnotesize 推荐流、帖子详情}

        \vspace{0.18cm}
        {\color{Sage}\footnotesize\bfseries 可执行动作}\par
        {\footnotesize 浏览、评论、转发、发帖}

        \vspace{0.18cm}
        {\color{Sage}\footnotesize\bfseries 可复盘结果}\par
        {\footnotesize 曝光、互动、记忆线索}
      \end{column}
      \begin{column}{0.54\paperwidth}
        \centering
        \vspace{0.08cm}
        \includegraphics[width=\linewidth,height=5.02cm,keepaspectratio]{slides/社交平台.png}
      \end{column}
    \end{columns}
\end{frame}

\begin{frame}{野营场景}
    \begin{columns}[T,totalwidth=0.92\paperwidth]
      \begin{column}{0.34\paperwidth}
        \small
        MetaSim 还有以下场景验证能力。

        \vspace{0.22cm}
        野营场景用于验证长期生存行为模拟。

        \vspace{0.28cm}
        {\color{Sage}\footnotesize\bfseries 环境压力}\par
        {\footnotesize 资源、体力、天气}

        \vspace{0.18cm}
        {\color{Sage}\footnotesize\bfseries 行为选择}\par
        {\footnotesize 采集、取水、修缮、休息}

        \vspace{0.18cm}
        {\color{Sage}\footnotesize\bfseries 长期轨迹}\par
        {\footnotesize 状态、背包、任务、记忆}
      \end{column}
      \begin{column}{0.54\paperwidth}
        \centering
        \vspace{0.08cm}
        \includegraphics[width=\linewidth,height=5.02cm,keepaspectratio]{slides/野营_compressed.png}
      \end{column}
    \end{columns}
\end{frame}

\begin{frame}{幼儿园陪伴场景}
    \begin{columns}[T,totalwidth=0.92\paperwidth]
      \begin{column}{0.34\paperwidth}
        \small
        MetaSim 还有以下场景验证能力。

        \vspace{0.22cm}
        幼儿园场景用于验证虚拟生活与现实陪伴的闭环。

        \vspace{0.28cm}
        {\color{Sage}\footnotesize\bfseries 虚拟生活}\par
        {\footnotesize 日程、通知、邀请、消息}

        \vspace{0.18cm}
        {\color{Sage}\footnotesize\bfseries 现实交互}\par
        {\footnotesize 机器人承载自然对话}

        \vspace{0.18cm}
        {\color{Sage}\footnotesize\bfseries 记忆回流}\par
        {\footnotesize 用户回应影响后续行动}
      \end{column}
      \begin{column}{0.54\paperwidth}
        \centering
        \vspace{0.08cm}
        \includegraphics[width=\linewidth,height=5.02cm,keepaspectratio]{slides/幼儿园_compressed.png}
      \end{column}
    \end{columns}
\end{frame}

\end{document}


